\section{Limitations}
\label{sec:limitations}

\subsection{Internal validity}
\label{sec:lim-internal}

The detection and false-flag rates are properties of the generators of histories, and the generators encode the assumptions of the experiments. The planted change is a persistent lengthening added at each breach date with probability $\pi_m$ and at least once, of 0.5 to 2 years on a stated parameter and of 0.5 to 4 years on a fitted one, and the unrelated changes of a stated parameter are independent across dates with a common probability. A bank whose unrelated changes cluster at the dates of the governance calendar, or whose override is made in several small steps, would give other rates; Proposition~\ref{prop:false-flag} covers the first case only with a $\lambda_t$ that varies, and the experiments do not. The breach chain has one persistence, and the closed forms and the measured rates depend on it through the number of breach dates. The comparators were run at their conventional settings: the Mann--Whitney test uses its asymptotic form, which is approximate for histories with many ties and for a dozen dates, and its measured size of \fitNullAMw{} to \fitNullCMw{} on the re-estimated population is above the nominal 5\,\%, whereas the Welch and Fisher tests are conservative or close to nominal. The change-point methods of Section~\ref{sec:related-work} are not included as separate detectors, because they have no input for the breach flags; the breach-agnostic count stands for them, and a method that estimated a break at an unknown date and then asked whether it fell on a breach would be a different test. The areas under the ROC curve are estimated on one set of 2,000 banks with a bootstrap interval, and the rates on five blocks whose standard deviations are in the committed results.

\subsection{External validity}
\label{sec:lim-external}

The committed history is one bank at three dates, and its behaviour is synthetic: no result here says what the bank did, and none says that a bank that follows the supervisory finding would be flagged. The simulated banks are block bootstraps of the 40 real quarterly rows of one bank in one regime, a decade of low and then rapidly rising rates, with blocks of six quarters, and the comparison with the real fitted series shows that they have more tiny and more large changes than the real series, which sits at the 95th percentile of the simulated median change. The calibrated tolerance inherits that mismatch: it holds its size on the simulated banks and is conservative on the real series. The block length was fixed before any detection experiment and was not tuned to the detection results; another length would move the tolerance. The histories are twelve quarters long (twenty-five on the real series, of which twelve are used), the number of breach dates is of the order of four, and a validator with a longer history or a different number of breach dates would face other rates; the closed forms show how, and the calibration has to be redone. The replicating portfolio of seven tranches with a 16-quarter window is the estimator of the product, and a bank with a different estimator, a deposit beta fitted on a different window, or a stochastic model of the deposit, would have a different distribution of changes and, perhaps, no ties at all.

\subsection{Construct validity and scope}
\label{sec:lim-construct}

A coincidence is not intent. The test measures the coincidence of a change with a flag, and the flag is a model output, the outlier ratio of a standardised shock, and not a limit that a committee set and recorded; a bank whose internal limit framework differs from the standard would give other flags. The dating convention, in which the flag of a date concerns the window ending there, is a choice, and an assumption changed in anticipation of a breach is by definition not a coincidence. The test sees only the parameters that are supplied to it: an assumption that is changed together with another, or replaced by a different parametrisation, is seen only if the series supplied reflects it. The internal gate that turns a count into a red light is the product's own and is not a regulatory limit. The tolerance family and its calibration are evaluated here and are not in the engine, so no claim of this paper concerns a released implementation of them. Finally, the absence claim of the introduction, that no published test of this kind was found, is bounded by the documents read; it is not a statement that none exists.
